\documentclass{rlreport}
\usepackage{listings}
\usepackage{xcolor}

% Define custom colors for Python syntax highlighting
\definecolor{codegreen}{rgb}{0,0.6,0}
\definecolor{codegray}{rgb}{0.5,0.5,0.5}
\definecolor{codepurple}{rgb}{0.58,0,0.82}
\definecolor{backcolour}{rgb}{0.95,0.95,0.92}

\lstdefinestyle{mystyle}{
    backgroundcolor=\color{backcolour},
    commentstyle=\color{codegreen},
    keywordstyle=\color{magenta},
    numberstyle=\tiny\color{codegray},
    stringstyle=\color{codepurple},
    basicstyle=\ttfamily\footnotesize,
    breakatwhitespace=false,
    breaklines=true,
    captionpos=b,
    keepspaces=true,
    numbers=left,
    numbersep=5pt,
    showspaces=false,
    showstringspaces=false,
    showtabs=false,
    tabsize=2
}

\lstset{style=mystyle}

\begin{document}

\title{RL-Based Adaptive Battery Management for UAVs}

\author{Laxmi \and Yash Raj \and Deva \and Nikita \and Ujwal}

\maketitle

\begin{abstract}
This is the abstract placeholder. 
\end{abstract}

\begin{keywords}
reinforcement learning, UAV, battery management, DQN, PPO, SAC, TD3
\end{keywords}

\section{Introduction}
This is the abstract placeholder.

 

\section{Methodology}

\subsection{Dataset Description (If any)}
No external dataset is used. All training data is generated by interacting with the simulated UAV environment.

\subsection{Problem Formulation: Markov Decision Processes}
\label{sec:mdp}
We formulate the problem as a Markov Decision Process (MDP), defined by the tuple $(\mathcal{S}, \mathcal{A}, P, R, \gamma)$. The agent maximizes the expected discounted return
\begin{equation}
  G_t = \sum_{k=0}^{\infty} \gamma^{k} R_{t+k+1}, \qquad 0 \le \gamma < 1.
  \label{eq:return}
\end{equation}

\subsubsection{States}
States description will be written here.

\subsubsection{Actions}
Actions description will be written here.

\subsubsection{Transition dynamics}
Transition dynamics will be written here.

\subsubsection{Episode termination}
Episode termination will be written here.

\subsection{Formulate reward}
Reward function will be written here.

\subsection{Establish value functions (in terms of your environmental variables)}
The action-value function satisfies the Bellman optimality equation
\begin{equation}
  Q^{*}(s,a) = \mathbb{E}\!\left[ R_{t+1} + \gamma \max_{a'} Q^{*}(S_{t+1}, a') \,\middle|\, S_t = s, A_t = a \right].
  \label{eq:bellman}
\end{equation}
Equation~\eqref{eq:bellman} follows from \eqref{eq:return}; see Section~\ref{sec:mdp}.

\subsection{RL Algorithms used in your project}
RL algorithms will be written here.
\subsection{Experimental Setup}
\label{sec:setup}

\subsubsection{Environment}
We use the frozen environment \texttt{env-v1}, which integrates the rotary-wing power model, the battery model (SoC, Peukert effect, 1-RC Thevenin, thermal, degradation) and the reward function.

\begin{table}[t]
\centering
\small
\caption{Environment parameters.}
\label{tab:env_params}
\begin{tabular}{lll}
\hline
Parameter & Value & Unit \\
\hline
UAV mass & TODO & kg \\
Payload & TODO & kg \\
Battery capacity & TODO & Wh \\
Wind range & TODO & m/s \\
Discount factor $\gamma$ & 0.99 & -- \\
\hline
\end{tabular}
\end{table}

\subsubsection{Training Protocol}
All agents are implemented with Stable-Baselines3 \cite{raffin2021sb3} on a Gymnasium \cite{towers2024gymnasium} environment. Each algorithm is trained with seeds 0, 1 and 2, and training is logged with the SB3 \texttt{Monitor} wrapper.

\subsubsection{Evaluation Protocol}
Every method is evaluated with the common script \texttt{evaluate.py} over 100 episodes with seeds 1000--1099. Results are reported as mean $\pm$ standard deviation over the three training seeds.

\subsubsection{Hyperparameters}
\begin{table}[t]
\centering
\small
\caption{Hyperparameters per algorithm.}
\label{tab:hyperparams}
\resizebox{\columnwidth}{!}{%
\begin{tabular}{lcccc}
\hline
Parameter & DQN & PPO & SAC & TD3 \\
\hline
Total steps & 300K & 500K & TODO & TODO \\
Learning rate & TODO & TODO & TODO & TODO \\
Batch size & TODO & TODO & TODO & TODO \\
Network size & TODO & TODO & TODO & TODO \\
$\gamma$ & 0.99 & 0.99 & 0.99 & 0.99 \\
\hline
\end{tabular}}
\end{table} 


\section{Results}
\label{sec:results}
This is the results placeholder. 

\section{Conclusion}
This is the conclusion placeholder.

% ---- References (stored separately in ref.bib) ----
\bibliographystyle{ieeetr}
\bibliography{ref}

\section*{Appendix}
\subsection{Code}
Code for the project goes here.

\subsection{Anything else}
Additional supporting figures or tables.

\end{document}